\documentclass[10pt,a4paper]{article}
\usepackage{fontspec,amsmath,amssymb,graphicx}
\usepackage[margin=18mm]{geometry}
\setmainfont{Times New Roman}
\setlength{\parindent}{0pt}
\setlength{\parskip}{5pt}
\allowdisplaybreaks
\sloppy
\emergencystretch=2em
\begin{document}
\section*{BT\_3\_6\_3\_03 — Lời giải bổ sung}
\subsection*{Kiến thức và trình tự giải}
\textbf{Bước 1 — Chọn Kirchhoff (3D) hoặc Poisson (2D).} Với li độ đầu bằng không, vận tốc $\psi$, tốc độ $c$, đặt $\rho=ct$. Trong ba chiều,


\[
u(P,t)=\frac t{4\pi}\int_{S^2}\psi(P+\rho\omega)d\omega.
\]


Trong hai chiều,


\[
u(P,t)=\frac1{2\pi c}\int_{|Q-P|<ct}\frac{\psi(Q)}{\sqrt{c^2t^2-|Q-P|^2}}dQ.
\]


Với li độ $f$ khác không, thêm đạo hàm theo $t$ của chính biểu thức tương ứng khi thay $\psi$ bằng $f$. Với nguồn và dữ liệu đầu không, tích phân theo thời gian của nghiệm có vận tốc là nguồn tại từng thời điểm (Duhamel).

\textbf{Bước 2 — Tận dụng hình học và đối xứng.} Chỉ thị miền $D$ biến Kirchhoff thành $t$ nhân tỷ lệ diện tích mặt cầu nằm trong $D$. Chọn trục qua tâm miền và $P$, đặt $\mu=\cos\theta$, $d\omega=d\mu\,d\phi$. Nếu điều kiện không phụ thuộc $\phi$, tỷ lệ là $|E|/2$ với $E\subset[-1,1]$. Mỗi nửa không gian cho một bất đẳng thức tuyến tính theo $\mu$; quả cầu cho một ngưỡng chỏm. Miền hợp/rút bỏ cho tổng/hiệu nếu các phần không chồng diện tích.

\textbf{Bước 3 — Viết chính xác các khoảng hợp lệ.} Giải bất đẳng thức gốc, sau đó mới bình phương khi có căn. Chia $[-1,1]$ ở nghiệm phương trình biên, kiểm tra dấu ở trung điểm từng khoảng, loại nghiệm ngoại lai. Nếu còn góc phương vị, dùng độ dài tập $\cos\phi<v$ rồi còn một tích phân theo $\mu$. Với Poisson, tích phân theo một tọa độ cho hiệu arcsin; cắt các cận vào đĩa trước khi tính.

\textbf{Bước 4 — Tính nguồn bằng nguyên hàm thời gian.} Với quả cầu bán kính $R$ cách $P$ khoảng $r>0$, đặt


\[
K_R(r,t)=t\left[\frac{R^2-(r-ct)^2}{4rct}\right]_0^1,\quad J_R(r,t)=\int_0^tK_R(r,s)ds.
\]


Tách ba giai đoạn: cả mặt cầu nằm trong miền, giao một chỏm, không giao. Nguồn bật trên $[0,T]$ cho $J_R(r,t)-J_R(r,(t-T)_+)$. Trường hợp $r=0$ tính trực tiếp, tránh chia cho $r$.

\textbf{Bước 5 — Kiểm tra chuẩn hóa, mặt sóng và cực đại.} Nếu $\psi=1$ trên toàn không gian thì $u=t$, dùng để kiểm tra hệ số. Với dữ liệu bước nhảy, chỉ kết luận cổ điển ngoài các mặt sóng. Cực đại phải xét cả điểm chuyển nhánh và giá trị một phía; nêu rõ supremum nếu chỉ đạt theo giới hạn. Nếu dùng tích phân số, ghi tham số và cách kiểm tra hội tụ; đồ thị không tự chứng minh một cực đại toàn cục.
\subsection*{Lời giải áp dụng cho bài}
\textbf{Bước 1. Đưa điều kiện xuyến về một biến góc.} Trên trục $z$, bán kính trụ của điểm cầu là $\rho\sqrt{1-\mu^2}$, chiều cao $z+\rho\mu$. Thay vào $(r'-R)^2+z'^2<a^2$, khai triển để được tập $E$ dưới đây. Bình phương phương trình biên có thể sinh nghiệm ngoại lai, nên chia các khoảng và kiểm tra bất đẳng thức gốc.

\textbf{Bước 2. Ở mặt phẳng $z=0$, tính góc phương vị.} Chọn $P=(r,0,0)$, đặt $h=\sqrt{a^2-\rho^2\mu^2}$. Điều kiện xuyến tương đương $R-h<r'<R+h$, với $|\rho\mu|<a$. Vì $R>a$, cận dưới dương nên có thể bình phương hai vế. Thay $r'^2=r^2+\rho^2(1-\mu^2)+2r\rho\sqrt{1-\mu^2}\cos\phi$ cho $C_-<\cos\phi<C_+$. Độ dài góc là hiệu $A(C_+)-A(C_-)$.

\textbf{Bước 3. Tích phân một chiều và cực trị.}

(a) Đặt $\rho=ct$. Tại $(0,0,z)$, điều kiện nằm trong xuyến là
\[E=\{\mu\in[-1,1]:\rho^2+R^2+z^2+2z\rho\mu-2R\rho\sqrt{1-\mu^2}<a^2\}.\]
Do đó $u=t|E|/2$. Tính các đầu mút bằng nghiệm của
\[(\rho^2+R^2+z^2-a^2+2z\rho\mu)^2=4R^2\rho^2(1-\mu^2),\]
kiểm tra bất đẳng thức gốc ở từng khoảng để loại nghiệm ngoại lai. Đây là phương trình bậc hai theo $\mu$, nên kết quả hữu hạn phép toán. Với $z=0$, miền đúng là $|\mu|<\sqrt{1-h^2}$ khi $h=(\rho^2+R^2-a^2)/(2R\rho)<1$, nên $u=t\sqrt{1-h^2}$; ngoài $R-a<\rho<R+a$, $u=0$.

(b) Với $r=\sqrt{x^2+y^2}>0$, $|\mu|<\min(1,a/\rho)$, đặt $h=\sqrt{a^2-\rho^2\mu^2}$, $D=2r\rho\sqrt{1-\mu^2}$,
\[C_\pm=\frac{(R\pm h)^2-r^2-\rho^2(1-\mu^2)}D.\]
Độ dài góc phương vị là $A(C_+)-A(C_-)$ với $A$ định nghĩa ở bài nón. Vậy
\[U_r(t)=\frac t{4\pi}\int_{-\min(1,a/\rho)}^{\min(1,a/\rho)}[A(C_+)-A(C_-)]d\mu.\]
Chỉ cần xét $0<t<(r+R+a)/c$; tính các cực trị $U'=0$ và điểm chuyển nhánh, lấy lớn nhất. Mã Python kèm theo thực hiện bước tích phân/tối ưu với tham số tùy chọn. Riêng $r=0$, từ câu (a), $U^2=[4R^2\rho^2-(\rho^2+R^2-a^2)^2]/(4R^2c^2)$; cực đại tại $\rho=\sqrt{R^2+a^2}$, cho $\max U=a/c$.

\textbf{Bước 4. Đối chiếu và kết luận.} Ở $r=0$, $U^2$ là một đa thức bậc hai theo $\rho^2$ với hệ số đầu âm. Đạo hàm bằng $0$ tại $\rho^2=R^2+a^2$, nằm giữa $(R-a)^2,(R+a)^2$, nên đây là cực đại toàn cục $a/c$. Ở $r>0$, biểu thức tích phân và cực trị một biến là chính xác theo tham số; kết quả từ mã là xấp xỉ, cần tăng cả độ chính xác tích phân và số khoảng thời gian để kiểm tra ổn định.

\textbf{Chi tiết nghiệm bậc hai tại điểm $(0,0,z)$.} Đặt $A=\rho^2+R^2+z^2-a^2$. Phương trình biên sau bình phương là
\[4\rho^2(z^2+R^2)\mu^2+4Az\rho\mu+A^2-4R^2\rho^2=0.\]
Với $\rho>0$, hai nghiệm thực nếu biểu thức dưới căn không âm:
\[\mu_\pm=\frac{-Az\pm R\sqrt{4\rho^2(z^2+R^2)-A^2}}{2\rho(z^2+R^2)}.\]
Giữ các nghiệm trong $[-1,1]$ và kiểm tra $A+2z\rho\mu=2R\rho\sqrt{1-\mu^2}$ để loại nghiệm ngoại lai. Chia khoảng ở các nghiệm còn lại rồi kiểm tra bất đẳng thức gốc trên từng khoảng để cộng độ dài; tại $t=0$ dùng dữ liệu $u=0$.

\textbf{Chi tiết khai triển biên và cực đại trên trục.} Ở (a), $(\rho\sqrt{1-\mu^2}-R)^2+(z+\rho\mu)^2$ khai triển thành $\rho^2+R^2+z^2+2z\rho\mu-2R\rho\sqrt{1-\mu^2}$: hai hạng $\rho^2\mu^2$ triệt tiêu. Với $z=r=0$, đặt $q=\rho^2$,
\[4R^2c^2U^2=4R^2q-(q+R^2-a^2)^2.\]
Đạo hàm theo $q$ bằng $4R^2-2(q+R^2-a^2)=2(R^2+a^2-q)$, đạo hàm bậc hai bằng $-2$. Cực đại duy nhất là $q=R^2+a^2$; thay vào được $4R^2c^2U^2=4R^2a^2$, tức $U=a/c$. $q-(R-a)^2=2Ra>0$ và $(R+a)^2-q=2Ra>0$, nên điểm cực đại nằm trong miền có sóng. Ở (b), $C_-<\cos\phi<C_+$ cho độ dài phương vị bằng hiệu hai độ dài $A(C_+)-A(C_-)$ sau khi cắt cận.
\begin{center}\includegraphics[width=.97\linewidth]{Do_thi/cone_torus.png}\end{center}\end{document}